% Hafta 2 — Saldırı ağacı nasıl kurulur?
% Derleme: py -3.12 tools/latex/derle.py docs/week-2/assets/h02-12-saldiri-agaci-kurulum.tex
\documentclass[tikz,border=8pt]{standalone}
\usepackage{cen429-cizim}
\begin{document}
\begin{tikzpicture}[node distance=8mm and 8mm]
  \node[baslik] (b) at (0,0) {Saldırı ağacı nasıl kurulur?};
  \node[altbaslik] at (b.south west) {niteliksel korkudan niceliksel karara};

  \node[kotukutu=38mm, minimum height=18mm, below=13mm of b.south west, anchor=north west, xshift=8mm] (s1)
    {\kb{Hedefi yaz}\\saldırganın\\kazanç tanımı};
  \node[kutu=38mm, minimum height=18mm, right=of s1] (s2) {\kb{Yolları listele}\\her yol bir dal};
  \node[kutu=38mm, minimum height=18mm, right=of s2] (s3) {\kb{VE / VEYA koy}\\hepsi mi gerekli,\\biri yeter mi?};
  \draw[ok, draw=soluk] (s1) -- (s2);
  \draw[ok, draw=soluk] (s2) -- (s3);

  \node[morkutu=38mm, minimum height=18mm, below=16mm of s1] (s4)
    {\kb{Yaprakları puanla}\\maliyet, beceri,\\erişim};
  \node[uyarikutu=38mm, minimum height=18mm, right=of s4] (s5) {\kb{En ucuz yolu bul}\\saldırgan onu seçer};
  \node[iyikutu=38mm, minimum height=18mm, right=of s5] (s6) {\kb{Oraya önlem koy}\\maliyeti yükselt};
  \draw[ok, draw=soluk] (s4) -- (s5);
  \draw[ok, draw=soluk] (s5) -- (s6);
  \draw[ok, draw=soluk] (s3.south) to[out=-90, in=90] (s4.north);

  \node[dip, text width=175mm, below=10mm of s4.south -| s5, anchor=north]
    {Ağaç, ``neyi önce düzeltelim?'' sorusuna sayıyla cevap verir.};
\end{tikzpicture}
\end{document}
